% =========================================================
% Proof: Density of Q in R
% Source: volume-iii/analysis/bounding/notes/notes-completeness.tex
% Mode: standard
% =========================================================

\subsection*{Density of $\mathbb{Q}$ in $\mathbb{R}$}
\label{prf:density-of-rationals}

\begin{remark*}[Return]
$\leftarrow$ \hyperref[thm:density-of-rationals]{Back to Theorem in Notes}
\end{remark*}

\begin{proof}
\Claim For every $a < b$ in $\mathbb{R}$, there exists $q \in \mathbb{Q}$
with $a < q < b$.

\Strategy
Scale the interval $(a,b)$ by $n$ until it has width greater than $1$,
guaranteeing it contains an integer $m$.
Then $q = (m+1)/n$ is the required rational.
The width condition $n(b-a) > 1$ is achieved by the Archimedean Property;
the integer inside the scaled interval is found by the Floor Lemma.

\medskip
Since $b - a > 0$, by the \ByThm{Archimedean Property} there exists
$n \in \mathbb{N}$ such that
\[
  n > \frac{1}{b-a}.
\]
Rearranging:
\[
  nb - na > 1.
\]
By the \ByThm{Floor Lemma}, there exists $m \in \mathbb{Z}$ with
\[
  m \leq na < m+1.
\]
Set $q := \dfrac{m+1}{n}$.
Since $m+1 \in \mathbb{Z}$ and $n \in \mathbb{N}$, we have $q \in \mathbb{Q}$.

\medskip
\textbf{Lower bound.}
From $na < m+1$, dividing by $n > 0$:
\[
  a < \frac{m+1}{n} = q.
\]

\medskip
\textbf{Upper bound.}
From $m \leq na$ and $nb > na + 1$:
\[
  nb > na + 1 \geq m + 1,
\]
and dividing by $n > 0$:
\[
  q = \frac{m+1}{n} < b.
\]

\Hence $a < q < b$ with $q \in \mathbb{Q}$. \AsReq
\end{proof}

\begin{remark*}[Proof shape]
Scale-and-floor construction.
The key insight: scale $(a,b)$ by $n$ until the width $n(b-a) > 1$,
which guarantees the interval $(na,nb)$ contains an integer.
The Archimedean Property produces the required $n$;
the Floor Lemma finds the integer $m$ with $m \leq na < m+1$;
then $m+1 \in (na,nb)$ and $q = (m+1)/n \in (a,b)$.

\FlashProof{Archimedean gives $n>1/(b-a)$, so $nb-na>1$. Floor Lemma gives $m\leq na<m+1$. Set $q=(m+1)/n\in\mathbb{Q}$. Lower: $na<m+1$ gives $a<q$. Upper: $nb>na+1\geq m+1$ gives $q<b$.}
\FlashUses{Archimedean Property; Floor Lemma; $\mathbb{Q}$ closed under $m/n$ with $m\in\mathbb{Z}$, $n\in\mathbb{N}$.}
\end{remark*}

\begin{remark*}[Why this completes $\overline{\mathbb{Q}} = \mathbb{R}$]
This theorem is the engine behind $\overline{\mathbb{Q}} = \mathbb{R}$.
Given any $x \in \mathbb{R}$ and $\varepsilon > 0$, apply Density of $\mathbb{Q}$
to the interval $(x - \varepsilon, x + \varepsilon)$ to obtain
$q \in \mathbb{Q}$ with $|x - q| < \varepsilon$.
Hence every real is adherent to $\mathbb{Q}$, so $\mathbb{R} \subseteq \overline{\mathbb{Q}}$.
See \hyperref[prf:closure-rationals]{Proof: $\overline{\mathbb{Q}} = \mathbb{R}$}.
\end{remark*}

\begin{remark*}[Dependencies]
\begin{itemize}
  \item \hyperref[prf:archimedean-property]{Archimedean Property}
  \item \hyperref[prf:floor-lemma]{Floor Lemma}
  \item $\mathbb{Q}$ is closed under $m/n$ for $m \in \mathbb{Z}$, $n \in \mathbb{N}$
\end{itemize}
\end{remark*}
